\documentclass{article}
\usepackage[utf8]{inputenc}
\usepackage{amsfonts}
% \usepackage{yfonts} % WIP
\usepackage{amsthm}
\usepackage{amsmath}
\usepackage{enumerate}
\usepackage{hyperref}
\usepackage{amssymb}
\usepackage{tikz} % diagram
\usepackage{changepage} % for adjustwidth

\begin{filecontents}[overwrite]{lattices.bib}
@misc{chris-peikert,
  author = {Chris Peikert},
  title = {{A decade of lattice cryptography}},
  year = {2016}
}
@misc{slides-micciano,
  author = {Daniele Micciancio},
  title = {{Introduction to Lattices}},
  year = {2024}
}
@misc{micciano-regev,
  author = {Daniele Micciancio and Oded Regev},
  title = {{Lattice-based Cryptography}},
  year = {2008}
}
@misc{matan,
  author = {Matan Prasma},
  title = {{Lattices, an elementary and rigurous account}},
  year = {2024}
}
@misc{menezes,
  author = {Alfred Menezes},
  title = {{A Gentle Introduction to Lattice-Based Cryptography}},
  year = {2026}
}
@misc{cynthia-dwork,
  author = {Cynthia Dwork},
  title = {{Lattices and Their Application to Cryptography}},
  year = {1998}
}
\end{filecontents}
\nocite{*}


\theoremstyle{definition}

\newtheorem{thm}{Theorem}[section]
\newtheorem{lemma}[thm]{Lemma}
\newtheorem{prop}[thm]{Proposition}
\newtheorem{cor}[thm]{Corollary}
\newtheorem{defn}[thm]{Definition}
\newtheorem{remark}[thm]{Remark}
\newtheorem{ex}[thm]{Example}

\newcommand{\Rr}{\mathbb{R}}
\newcommand{\Zz}{\mathbb{Z}}
\newcommand{\Ll}{\mathcal{L}}

\title{Notes on Lattices}
\author{arnaucube}
\date{2026}

\begin{document}

\maketitle

\begin{abstract}
    Notes taken while studying lattice-based cryptography.

    Usually while reading books and papers I take handwritten notes in paper notebooks, this document contains some of them re-written to $LaTeX$.
\end{abstract}

\tableofcontents

\section{Lattices}
\subsection{Hard problems}

\begin{defn}[SVP (Shortest Vector Problem)]
find a shortest non-zero vector $x \in \Ll$, ie. that minimizes the Euclidean norm $||x||$, that is $\lambda_1(\Ll)=||x||$.
\end{defn}

\begin{defn}[CVP (Closest Vector Problem)]
given a vector $y \in \Rr^n$ with $y \not\in \Ll$, find a nonzero vector $x \in \Ll$ that is closest to $y$, ie. that minimizes the Euclidean norm $||y-x||$.
\end{defn}

\begin{defn}[NTRU key recovery problem]
Given $h(x)$, find ternary polynomials $f, g \in \mathcal{T}$ satisfying
$$f(x) \cdot h(x) = g(x) \pmod q \in R_q$$
Where $R_q = \Zz_q[X]/(X^n-1)$.
\end{defn}

\begin{defn}[SIS (Short Integer Solution)]
Given $q \in \Zz$ and matrix $A \in \Zz_q^{m \times n}$, find a non-zero short $z \in \Zz^n$ s.th.
$$A \cdot z = 0 \pmod q$$

\small{"short": $z \in \{B, B\}^n$ for some bound $B$, often $z \in \{-1, 0, 1 \}^n$.}
\end{defn}

The set of all solutions to the SIS problem forms the (q-ary) lattice
$$\Lambda^{\perp}_q(A) = \{ z \in \Zz^n ~|~ A \cdot z = 0 \pmod q \}$$
where the SIS problem becomes the SVP for $\Lambda_q^{\perp}(A)$

\begin{itemize}
\item LLL algorithm is used to solve SVP and CVP.

  \item Both SVP and CVP become computationally difficult (NP hard*) as the lattice dimension $n$ grows.
\end{itemize}


\begin{defn}[Ring SIS (Ring Short Integer Solution)]
given a uniformly random matrix $A \in^R R_q^n$, with elements $a_i \in R_q$, find a non-zero integer vector $z \in R^n$ of norm $||z|| \leq \beta$ such that

$$A z = \sum_i a_i \cdot z_i = 0 \in R_q$$
\end{defn}

\begin{defn}[MSIS (Module Short Integer Solution)]
given a uniformly random matrix $A \in^R R_q^{m \times n}$, with columns $a_i \in R_q^m$, find a non-zero integer vector $z \in R^n$ of norm $||z|| \leq \beta$ such that

$$A z = 0 \in R_q^m$$
$$ie.~ \sum_i a_{ij} \cdot z_j = 0 \in R_q, ~~for~ i=1, \ldots, m$$
\end{defn}

Notice that Module-SIS is a Ring-SIS with $rank=1$ ($m=1$).

\subsection{Dual lattices}

From previous section, we have LWE-Lattice and SIS-Lattice:
\begin{defn}[LWE Lattice] \label{lwe-lattice}
  A LWE-Lattice $\Ll$ is defined
  \begin{align*}
    \Ll &= \{ A^T s ~|~ s \in \Zz_q^n \} + q \Zz^m\\
	&= \{ A^T s + qz ~|~ s \in \Zz_q^n, z \in \Zz^m \}
  \end{align*}
\end{defn}

\begin{defn}[SIS-Lattice] \label{sis-lattice}
  A SIS-Lattice $\Ll$ is defined
$$\Ll^\perp = \{ z \in \Zz^m ~|~ Az=0 \in \Zz_q^n \} \supseteq q \Zz^m$$
\end{defn}


\begin{defn}[Dual of a lattice] \label{dual}
    For a lattice $\Ll$, its dual is
    $$\Ll^* = \{ y \in \Rr^m ~|~ \langle y, x \rangle \in \Zz ~\forall~ x \in \Ll \}$$
\end{defn}

\begin{lemma}
    $$\Ll = q \Ll^{\perp *} ~\text{ie.}~ \frac{1}{q} \Ll = \Ll^{\perp *}$$
\end{lemma}
\begin{proof}
    \begin{enumerate}[i.]
    \item
    Take $y \in \frac{1}{q} \Ll$, then $\exists s \in \Zz^n, z \in \Zz^m$ such that
    $$y = \frac{1}{q} (A^Ts + qz) = \frac{1}{q} A^T s +z$$

    Let $x \in \Ll^\perp$. By definition (\ref{sis-lattice}) $Ax=0 \pmod q$, therefore $\exists~ k \in \Zz^n$ such that $Ax = qk$.

    Compute
    \begin{align*}
	\langle y, x \rangle &= \langle \frac{1}{q} A^T s + z, x \rangle\\
			     &= \frac{1}{q} \langle A^T s, x \rangle + \langle z, x \rangle\\
			     &= \frac{1}{q} \langle s, Ax \rangle + \langle z, x \rangle\\
			     &= \frac{1}{q} \langle s, qk \rangle + \langle z, x \rangle\\
			     &= \langle s, k \rangle + \langle z, x \rangle
    \end{align*}
    and $\langle s, k \rangle + \langle z, x \rangle \in \Zz$, since each of $s, k, z, x \in \Zz$.

    Thus $\langle y, x \rangle \in \Zz$, hence by \ref{dual}, $y \in \Ll^{\perp *}$.

    We started from $y \in \frac{1}{q}\Ll$, so, now we have
    $$y \in \Ll^{\perp *} ~\forall~ y \in \frac{1}{q} \Ll$$
    therefore
    $$\Ll^{\perp *} \supseteq \frac{1}{q} \Ll$$
    
    \item
    Now we prove the reverse inclusion. Take $y \in \Ll^{\perp *}$.

    Observe that $q \Zz \subseteq \Ll^\perp$, because
    $$\forall z \in \Zz^m,~ A(qz) = qAz = q 0 = 0 \pmod q$$

    Since $y \in \Ll^{\perp *}$, we must have $\langle y, qz \rangle \in \Zz ~\forall~ z \in \Zz^m$.

    Let $e_i \in \Zz^m$ be the $i-th$ standard basis vector, ie. $e_i=(0, \ldots, 0, 1, 0, \ldots, 0)$ (where one is at the i-th position).

    We know that $\langle y, qz \rangle \in \Zz ~\forall z \in \Zz^m$ (including $z=e_i$), thus\\
    $q e_i = (0, \ldots, 0, q, 0, \ldots, 0)$. Therefore,
    $$\langle y, q e_i \rangle = q \langle y, e_i \rangle = q \cdot y_i$$
    since $\langle y, e_i \rangle = y_i$ (extracts the i-th position from $y$).

    So, since $\langle y, qz \rangle \in \Zz ~\forall z \in \Zz^m$, and $\langle y, q e_i \rangle = q y_i$, then $q y_i \in \Zz$

    This argument works $\forall i \in [m]$, that is, $q y_1, \ldots, q y_m \in \Zz$. Therefore
    $$qy \in \Zz^m ~\forall~ y \in \Ll^{\perp *}$$

    $qy \pmod q$ is orthogonal to $ker(A)$.  %TODO

    Let $\overline{x} \in ker(A) \subseteq \Zz_q^m$.
    Since $A \overline{x}=0$, then $Ax =0 \pmod q$, thus $x \in \Ll^\perp$.

    Since $y \in \Ll^{\perp *}$, by the definition of $\Ll^{\perp *}$, $\langle y, x \rangle \in \Zz$. Multiply it by $q$,
    $$ q \langle y, x \rangle = \langle qy, x \rangle \in q \Zz$$
    hence $\langle qy, x \rangle = 0 \pmod q$, thus $qy \pmod q \in ker(A)^\perp$.

    \begin{adjustwidth}{4em}{0pt}
	Observe: $ker(A)^\perp = im(A^T)$, proof:
	let $A: \Zz_q^m \longrightarrow \Zz_q^n$, define
	$$ker(A) = \{ x \in \Zz_q^m ~|~ Ax=0 \} = \Ll^\perp(A)$$

	Define the orthogonal complement
	$$(ker(A))^\perp = \{ y \in \Zz_q^m ~|~ \langle y, x \rangle =0 \pmod q ~\forall~ x \in ker(A) \}$$
	$$im(A^T) = \{ A^T s ~|~ s \in \Zz_q^n \}$$

	Then, $(ker(A))^\perp = im(A^T)$.
    \end{adjustwidth}

    Continuing from $qy \pmod q \in (ker(A))^\perp$, we've just seen that $(ker(A))^\perp = im(A^T)$, thus
    $$\exists s \in \Zz^n ~\text{s.th.}~ qy=A^Ts \pmod q$$

    Therefore, for some
    $$z \in \Zz^m, ~ qy=A^Ts +qz$$
    which means that $qy \in \Ll$, by definition of $\Ll$ (\ref{lwe-lattice}). Divide that by $q$, obtaining
    $$y \in \frac{1}{q} \Ll$$

    Thus, $\forall y \in \Ll^{\perp *}$, we have $y \in \frac{1}{q} \Ll$.

    Hence $\Ll^{\perp *} \subseteq \frac{1}{q} \Ll$.
    \end{enumerate}

    \vspace{0.6cm}
    Putting together $(i).~ \Ll^{\perp *} \supseteq \frac{1}{q} \Ll$, and $(ii).~ \Ll^{\perp *} \subseteq \frac{1}{q} \Ll$, we have
    $$\Ll^{\perp *} = \frac{1}{q} \Ll$$
\end{proof}


\bibliographystyle{unsrt}
\bibliography{lattices.bib}

\end{document}
